% What this holds: Section Election Method.
% Read by arlington-bsap.tex. Prose lives here; formatting lives in the wrapper and style/.
\section{Election Method}\label{sec:method}

Arlington has elected its Board three ways since 1870, and what separates them is
how many seats one ballot fills and who counts as the electorate for each.

\subsubsection{Districts, At-Large Election and Ranked Choice}

Under single-winner districts the county is divided into areas and each elects one
member, so a voter has one vote in one seat. A district favors a minority
concentrated in one neighborhood: the Supreme Court asks whether such a group is
large and geographically compact enough to make a majority in a single-member
district, and where its members are scattered they cannot.\autocite[50]{thornburgvgingles1986}

Election at large puts every seat to every voter in the county. When one seat is
open on the ballot, a voter has one vote, as in a district, but the electorate is
the whole county, so a group must carry the county and its concentration in one
neighborhood no longer counts. When several seats are open on one ballot, the
method is block voting: every voter votes for as many candidates as there are seats,
and the candidates with the most votes win.\autocite{arlingtonboard2026rcv} A block vote
favors the largest group. If the voters in the biggest bloc support the same
candidates, that bloc picks every winner and the rest of the community picks none.
In Lynchburg in 2022 three Republican nominees took all three at-large council
seats, with 12,287, 11,658 and 10,685 votes, while the leading independent finished
528 votes behind the third of them.\autocite[contest 156756]{vaelections}\fnsep\autocite[48--49]{thornburgvgingles1986}
That is what makes a block vote hyper-majoritarian: the rule lets a bare majority
voting as one take every seat on the ballot and leaves a smaller bloc none.

In the account of a 1979 University of Virginia study, block voting spread because
the reformers who adopted it saw it as a weapon: in many cities in the early 1900s,
``upper-class, business-led reformers regarded at-large elections as a device to
reduce the influence of working class and ethnic neighborhoods that found
expression through ward based elections.''\autocite{orourke1979atlarge}
A study of 1976 council seats found what the device could do: Black residents of
cities electing at large held about half the representation their share of the
population would warrant, slightly less than half of what Black residents of ward
cities held.\autocite[122]{berry1979atlarge}

Under ranked choice voting, which the General Assembly calls instant runoff voting,
a voter ranks candidates in order of preference and the count proceeds in rounds,
either electing a candidate or defeating the last-place candidate and passing those
votes to the voters' next choices, until the seats are
filled.\autocite{vaacts2020c713} In a single-winner count, a candidate needs more
than half the votes, and the candidate with the fewest is eliminated round by round
until one passes half.\autocite{arlingtonboard2026rcv} In a multi-winner count, the
single transferable vote, a candidate needs a smaller share set by the number of
seats, and a winner's surplus passes to the voters' next
choices.\autocite{vaacts2020c713}
Arlington used the multi-winner count once, in the Democratic primary of 2023.

\subsubsection{Districts, 1870--1930}

For its first 60 years the Board was chosen district by district. Arlington,
Jefferson and Washington each elected one supervisor, so a voter had one vote and
a say in one seat of three. Through 1901 a supervisor served a single year, from
July 1 after the May election, and his district's voters were back at the polls the
next spring.\autocite[art.~VII, sec.~2]{vaconstitution1869}\fnsep\autocite[ch.~76,
secs.~4 and 14]{vaacts1870} The constitution of 1902 moved county elections to
November and lengthened the term to four years beginning each January 1; the
supervisors elected in May 1901 served 30 months, to January 1, 1904, and the first
November winners took their seats
then.\autocite[sec.~112; Schedule, sec.~10]{vaconstitution1902}
% waits on: nominations-pre-1931

\subsubsection{At Large, 1930 to the Present}

In November 1930 the county's voters ended the district system. In the referendum
that adopted the County Manager Plan they chose election at large over election by
district, 1,689 to 1,149.\autocite[194--196]{rose1976}
% waits on: method-1930-advocacy
The county's case for one electorate ran through \textit{Bennett v. Garrett}
(1922), in which the Supreme Court of Appeals had refused Clarendon a town charter
because the ``general good of the community'' the statute protects ``must be
construed as embracing the whole section'';\autocite{bennettvgarrett1922} Cornelia
Rose's county history reads that reasoning into the at-large method the county
chose in 1930.\autocite[176, 197]{rose1976}
% Rose's phrase "continuous, contiguous, and homogeneous community" is hers, not
% the court's: the bennettvgarrett1922 annotation and the county-corrections row.

The first at-large election, in November 1931, put 51 candidates before every voter
for five seats, and the whole Board stood at once.\autocite[67]{anderson1958}
% source: the 1931 ballot, to say how many votes each voter had for the five seats
That was a block vote in full, and the Board stood together again every four years
until 1938.

\subsubsection{Staggered Terms, 1938 to the Present}

In 1938 voters approved staggered terms, 1,539 to 1,487 in the county's
tally.\autocite[11]{arlingtonelections2021}\fnsep\autocite{sun1938referendum} The
ballot asked only whether members should be elected ``as provided in Sec.~2773-F1
of the Code of Virginia,'' and the Commonwealth's Attorney's office fielded calls
from voters asking what that meant. The Arlington County Woman's Democratic Club
opposed the plan, warning that it would make it ``virtually impossible for a woman
to be elected to the board.''\autocite{sun1938referendum}

Staggering changed the ballot as much as the calendar. From the Board elected in
1939, one member was chosen in most years and two in the fourth, in 1943, 1947,
1951 and every fourth year after.\autocite[37]{arlhist1967officials}
% Two-seat years are elections_turnout.csv board_seats = 2; the table counts 1939
% itself as a five-seat election.
In a two-seat year a voter cast two votes, as the county's returns for 1955 head
the County Board contest ``Vote for 2,'' and the two candidates with the most votes
won.\autocite[28]{arlingtonelections2021} That is block voting: each voter may
support as many candidates as there are seats, and the highest totals win.

In the other years no ballot was a sweep, because no ballot held more than one seat.
Every one of those seats, though, went to the same countywide electorate, so a
majority bloc that held together won each seat as it came up, and over one
four-year cycle it held all five. Staggering made five seats look like five
separate choices spread across four years. For a majority that votes as one they
are a single choice made in installments, and its result is the one a block vote
would have produced in 1931.

Though every seat went to the countywide electorate, no two went to the same
voters. The seat was one race on a ballot headed by something else, a President, a
governor or only Congress, and what headed the ballot decided how many people came
out. In the presidential year 2012 the Board's vote equalled
\boardVoteTwoThousandTwelve{} percent of residents of voting age; in the governor's
year after it, \boardVoteTwoThousandThirteen{} percent
(Figure~\ref{fig:turnout-board}). In each of the \boardPairs{} presidential years from
\boardPairsFirst{} through \boardPairsLast{} that can be set beside the governor's
year after it, the presidential year drew the larger vote. Two members of one Board,
elected a year apart, were chosen by electorates of different size.

\begin{figure}
  \centering
  \refstepcounter{figure}\label{fig:turnout-board}
  {\raggedright\MakeUppercase{Figure \thefigure}\par}
  {\raggedright\bfseries Votes Cast per 100 Residents of Voting Age, 1932--2024\par}
  \medskip
  \includegraphics[width=\textwidth]{elections_turnout_board.pdf}
  \smallskip
  \figurenotes{Notes: Residents of voting age are men aged 21 and over through
  1916, all residents 21 and over from 1920, when women first vote, and all
  residents 18 and over from 1971. Each is the census count, with a straight line
  between censuses, and includes non-citizens and the people the laws of the
  day disfranchised. A dashed rule marks 1940, the first staggered election;
  before it the Board's five seats were filled all at once, in 1931, 1935 and
  1939. From 1940 the Board's vote is drawn only for years that filled one
  seat, shaded by what led the ballot. A year that filled more than one seat,
  or whose count is incomplete, is left blank, because its votes are not its
  voters.}{Source: Warrock-Richardson Almanack; Secretary of the Commonwealth;
  Virginia Department of Elections; \textit{Alexandria Gazette}; O'Leary
  (2010); Arlington County, \textit{Arlington Elections}; US decennial census,
  1870--2020.}
\end{figure}

\subsubsection{Ranked Choice, 2020 to the Present}

The General Assembly authorized ranked choice voting in 2020, letting a county
under the Manager Plan provide by ordinance for instant runoff voting in the
nomination or election of its Board.\autocite{vaacts2020c713} A single-winner
count still fills one seat from one countywide electorate, by a majority of the
ballots that reach the final round. The County Board adopted ranked choice voting for all later primaries in December
2023, so the primaries of 2024 and 2025 each filled one nomination by a
single-winner count, and in February 2025 it extended the method to the November 2025
general election, a contest for one seat.\autocite{arlingtonboard2026rcv}

The multi-winner count of 2023 is the exception. With two nominations open among
six candidates in the state-run Democratic primary, the count was a two-winner
single transferable vote: a nominee's surplus passed to the voters' next choices.
Natalie Roy had more first choices than Maureen Coffey, 5,607 to 5,421, and Coffey
was nominated beside Susan Cunningham all the
same.\autocite[contest 160326]{vaelections}\fnsep\autocite{arlnow2023coffeyprimary}
% waits on: margins-ranked-choice-rounds, for the 2023 round-by-round count
Roy led on first choices and Coffey was nominated on transfers.
% parked: bridge sentence into Part C, Election Method, when Part C is worked out

% Party by seat waits on the County.
% waits on: party-unlabelled
% waits on: party-before-1932

